% Exercise bank — Chapter 4
% The chapter is determined by the name of this file.

\begin{exo}[Integrating a differential form along a path]{integrale-forme-differentielle}
\keywords{differential form, exact differential, Schwarz criterion, line integral}

Consider the two differential forms
\[
\omega_1 = y\,\mathrm{d}x + x\,\mathrm{d}y,
\qquad
\omega_2 = y\,\mathrm{d}x,
\]
and three paths joining the origin $O(0,0)$ to the point $M(1,1)$: path
$\gamma_1$, consisting of the horizontal segment $O\to(1,0)$ followed by the
vertical segment $(1,0)\to M$; path $\gamma_2$, consisting of the vertical
segment $O\to(0,1)$ followed by the horizontal segment $(0,1)\to M$; and path
$\gamma_3$, the straight-line segment joining $O$ directly to $M$.

\begin{enumerate}
\item Show that $\omega_1$ is exact by finding a function $f$ such that $\omega_1 = \mathrm{d}f$.
\item Calculate $\int_{\gamma_1}\omega_1$ and $\int_{\gamma_2}\omega_1$ by parametrizing each segment. Recover the result without calculation.
\item Apply Schwarz's criterion to $\omega_2$. What can be concluded?
\item Calculate $\int_{\gamma_1}\omega_2$, $\int_{\gamma_2}\omega_2$, and $\int_{\gamma_3}\omega_2$. Comment on the results.
\end{enumerate}

\begin{indication}
On a horizontal segment $\mathrm{d}y = 0$, and on a vertical segment
$\mathrm{d}x = 0$: each integral reduces to an ordinary single integral.
\end{indication}

\begin{solution}
\textbf{Question 1.} The function $f(x,y) = xy$ works, since $\partial f/\partial x = y$ and $\partial f/\partial y = x$.

\textbf{Question 2.} For a curve parametrized by
$t\mapsto\bigl(x(t),y(t)\bigr)$, we have
\[
\int_\gamma\omega_1
= \int \left[y(t)x'(t)+x(t)y'(t)\right]\,\mathrm{d}t.
\]

Path $\gamma_1$ is the union of two segments. On the first,
$x(t)=t$ and $y(t)=0$, with $t\in[0,1]$; hence $\mathrm{d}x=\mathrm{d}t$
and $\mathrm{d}y=0$. On the second, $x(t)=1$ and $y(t)=t$, again with
$t\in[0,1]$; hence $\mathrm{d}x=0$ and $\mathrm{d}y=\mathrm{d}t$. Therefore,
\[
\begin{aligned}
I_1
= \int_{\gamma_1}\omega_1
&= \int_0^1\left(0\times 1+t\times 0\right)\,\mathrm{d}t
 + \int_0^1\left(t\times 0+1\times 1\right)\,\mathrm{d}t \\
&= 0+\int_0^1\mathrm{d}t = 1.
\end{aligned}
\]

For $\gamma_2$, the first segment is parametrized by $x(t)=0$, $y(t)=t$,
and the second by $x(t)=t$, $y(t)=1$, with $t\in[0,1]$ in both cases.
We therefore obtain
\[
\begin{aligned}
I_2
= \int_{\gamma_2}\omega_1
&= \int_0^1\left(t\times 0+0\times 1\right)\,\mathrm{d}t
 + \int_0^1\left(1\times 1+t\times 0\right)\,\mathrm{d}t \\
&= 0+\int_0^1\mathrm{d}t = 1.
\end{aligned}
\]
The two values agree. This could have been predicted without calculation:
because $\omega_1=\mathrm{d}f$ with $f(x,y)=xy$, its integral depends only on
the endpoints,
\[
\int_\gamma\omega_1=f(M)-f(O)=f(1,1)-f(0,0)=1.
\]

\textbf{Question 3.} Write $\omega_2=A(x,y)\,\mathrm{d}x+B(x,y)\,\mathrm{d}y$.
Here $A(x,y)=y$ and $B(x,y)=0$. If the form were exact, Schwarz's criterion
would give
\[
\frac{\partial A}{\partial y}=\frac{\partial B}{\partial x}.
\]
However, $\partial A/\partial y=1$, whereas $\partial B/\partial x=0$.
The mixed derivatives differ, so $\omega_2$ is not exact. Its integral may
therefore depend on the path followed between $O$ and $M$, as the following
calculation confirms.

\textbf{Question 4.} Since $\omega_2=y\,\mathrm{d}x$, only a change in $x$
at a nonzero value of $y$ can contribute to the integral. Using the previous
parametrizations, on $\gamma_1$ we find
\[
\begin{aligned}
J_1
=\int_{\gamma_1}\omega_2
&=\int_0^1 0\times 1\,\mathrm{d}t
  +\int_0^1 t\times 0\,\mathrm{d}t \\
&=0.
\end{aligned}
\]
On $\gamma_2$, the first segment satisfies $x(t)=0$, so $\mathrm{d}x=0$;
on the second, $x(t)=t$, $y(t)=1$, and $\mathrm{d}x=\mathrm{d}t$. Thus
\[
\begin{aligned}
J_2
=\int_{\gamma_2}\omega_2
&=\int_0^1 t\times 0\,\mathrm{d}t
  +\int_0^1 1\times 1\,\mathrm{d}t \\
&=1.
\end{aligned}
\]
Finally, the diagonal segment $\gamma_3$ is explicitly parametrized by
\[
x(t)=t,\qquad y(t)=t,\qquad t\in[0,1],
\]
so $\mathrm{d}x=\mathrm{d}t$ and $\mathrm{d}y=\mathrm{d}t$. Hence
\[
J_3
=\int_{\gamma_3}\omega_2
=\int_0^1 y(t)x'(t)\,\mathrm{d}t
=\int_0^1 t\,\mathrm{d}t
=\frac12.
\]
The three paths have the same endpoints but yield three different values:
$J_1=0$, $J_2=1$, and $J_3=1/2$. This is precisely the path dependence
characteristic of a non-exact differential form.
\end{solution}
\end{exo}

\begin{exo}[Same change in internal energy, three modes of transfer]{meme-delta-u-trois-transferts}
\keywords{first law, internal energy, heat, mechanical work, electrical work, calorimetry, state function}

A liquid, its calorimeter, and a small immersed resistor form a closed system
that is macroscopically at rest. The container is rigid, and the total heat
capacity of the system is assumed constant:
\[
C=2{,}00\ \mathrm{kJ\,K^{-1}}.
\]
The system is to be taken from equilibrium state $A$, at
$T_A=20{,}0\,^\circ\mathrm{C}$, to equilibrium state $B$, at
$T_B=25{,}0\,^\circ\mathrm{C}$. It is given that
\[
U(B)-U(A)=C(T_B-T_A).
\]
The changes in macroscopic kinetic and potential energy are zero. All losses
to the surroundings are neglected.

The following three protocols are carried out independently, starting from
the same initial state $A$:
\begin{itemize}
\item[Protocol 1] The system is placed in thermal contact with a hot reservoir,
with no work supplied to it.
\item[Protocol 2] The walls are thermally insulated. A paddle stirs the liquid
as a mass $m$ falls through a height $h=5{,}00$ m. The paddle and liquid end
at rest, and all the mechanical energy lost by the mass is transferred to the
system. Take $g=9{,}81\ \mathrm{m\,s^{-2}}$.
\item[Protocol 3] The system receives heat $Q_3=4{,}00$ kJ, and the remaining
energy is supplied electrically to the resistor. The electrical power received
is constant and equal to $\mathcal P=300$ W.
\end{itemize}

\begin{enumerate}
\item Calculate the change in internal energy $\Delta U=U(B)-U(A)$ common to all three protocols.
\item For each of the first two protocols, determine the heat $Q$ and work $W$ received by the system. In the second protocol, calculate the required mass $m$.
\item For the third protocol, calculate the electrical work received and the duration for which the resistor is powered.
\item The three transformations have the same initial and final states. Explain why their values of $Q$ and $W$ may nevertheless differ. Does the system contain ``heat'' or ``work'' in the final state $B$?
\end{enumerate}

\begin{indication}
Apply $\Delta U=Q+W$ using the banker's convention. For the paddle, the work
received equals the decrease $mgh$ in the potential energy of the mass.
Electrical energy crossing the boundary through the wires is counted as
electrical work.
\end{indication}

\begin{solution}
\textbf{Question 1.} The temperature difference is $T_B-T_A=5{,}0$ K.
Therefore,
\[
\Delta U=C(T_B-T_A)
=2{,}00\times5{,}0
=10{,}0\ \mathrm{kJ}.
\]
This value depends only on the two equilibrium states $A$ and $B$.

\textbf{Question 2.} In the first protocol, no work is received: $W_1=0$.
The first law therefore gives
\[
Q_1=\Delta U=10{,}0\ \mathrm{kJ}.
\]
The heat is positive because energy enters the system.

In the second protocol, the walls are thermally insulated, so $Q_2=0$.
The entire change in internal energy comes from the mechanical work of the
paddle:
\[
W_2=\Delta U=10{,}0\ \mathrm{kJ}.
\]
Since $W_2=mgh$, it follows that
\[
m=\frac{W_2}{gh}
=\frac{10{,}0\times10^3}{9{,}81\times5{,}00}
\simeq 204\ \mathrm{kg}.
\]
The large order of magnitude is a reminder that even a modest temperature
rise already corresponds to a substantial amount of mechanical energy.

\textbf{Question 3.} The first law requires
\[
W_3=\Delta U-Q_3=10{,}0-4{,}00=6{,}00\ \mathrm{kJ}.
\]
This work is positive: the electrical supply provides energy to the system.
With $W_3=\mathcal P\,\Delta t$,
\[
\Delta t=\frac{W_3}{\mathcal P}
=\frac{6{,}00\times10^3}{300}
=20{,}0\ \mathrm{s}.
\]

\textbf{Question 4.} The three paths give, respectively,
\[
(Q_1,W_1)=(10{,}0,0)\ \mathrm{kJ},
\qquad
(Q_2,W_2)=(0,10{,}0)\ \mathrm{kJ},
\]
and
\[
(Q_3,W_3)=(4{,}00,6{,}00)\ \mathrm{kJ}.
\]
In every case, the sum $Q+W$ is $10{,}0$ kJ and produces the same change
$\Delta U$. Internal energy is a state function, whereas heat and work
describe transfers during a transformation. In final state $B$, the system
has internal energy, but it contains neither ``heat'' nor ``work.''
\end{solution}
\end{exo}

\begin{exo}[Compression under constant external pressure]{compression-pression-exterieure-constante}
\seoready{true}
\keywords{pressure work, constant external pressure, compression, expansion, free expansion}

A gas is compressed from a volume $V_A$ to a volume $V_B<V_A$ under a
constant external pressure $P_0$.

\begin{enumerate}
\item Calculate the work received by the gas.
\item The gas returns from $V_B$ to $V_A$ against the same external pressure $P_0$. Calculate the work received and compare it with the compression work.
\item Repeat the expansion from $V_B$ to $V_A$ when the gas expands into a vacuum. Interpret the three signs obtained.
\end{enumerate}

\begin{indication}
Use $W=-\int P_{\mathrm{ext}}\,\mathrm{d}V$. The \emph{external} pressure must
be integrated, including when the transformation is not quasi-static.
\end{indication}

\begin{solution}
\textbf{Question 1.} The external pressure is constant, so
\[
W_{A\to B}=-P_0(V_B-V_A)=P_0(V_A-V_B)>0.
\]
The gas receives work during compression.

\textbf{Question 2.} For the return path,
\[
W_{B\to A}=-P_0(V_A-V_B)<0.
\]
The two works are opposites because the transformations are carried out
against the same constant external pressure.

\textbf{Question 3.} In a vacuum, $P_{\mathrm{ext}}=0$, so
\[
W_{B\to A}^{\mathrm{vide}}=0.
\]
Thus, expansion does not automatically imply negative work: an external
force must actually be displaced.
\end{solution}
\end{exo}
